We thank the referee for distinguishing the sound original-model lower bound and feasible reconstruction from the paper's insufficiently sharp positive frontier and incomplete computational interpretation. This revision addresses that distinction with a new exact tractability theorem, a bounded-path extension, an exact price-support criterion and a prospectively frozen study. It preserves all earlier mathematical results and their complete proofs; neither a negative empirical outcome nor a difficult regime has been removed. References below point to the current readers.

\section*{1. Scientific hierarchy and a positive complexity frontier}
\textbf{Concern.} The original-model W[1] reduction is useful, but hardness plus a pseudo-polynomial grid does not sufficiently identify when the model is tractable.

\textbf{Response.} Section~\ref{sec:tariff60} adds Theorem~\ref{thm:tariff60}. For common linear rewards and free service, expected caps, realized ceilings, probabilities, menu allowance and rational denominators remain unrestricted. A standard opening fee with $d$ exceptional commands admits exact optimization in $O(kN^2+2^d mN^2)$ rational operations and $2^d\operatorname{poly}(L)$ bit time. The exponent is independent of $d$. Uniform nonnegative fees give a polynomial algorithm with a strongly polynomial arithmetic count. Lemma~\ref{lem:capacity60} proves the capacity identity and recovers a feasible original policy, rather than solving only an auxiliary path problem. Conditioning on the installed exceptional subset makes charges constant at fixed book size; a layered longest-capacity recurrence is then sufficient. This positive result uses neither fixed menu size nor bounded denominators.

The title, abstract and introduction now organize the paper around this tariff-sensitive frontier. Corollary~\ref{cor:fee60} derives the exact uniform-fee value envelope and optimal-cardinality monotonicity. The prior W[1] reduction is retained in full. Uniform fees with general curved rewards and costly service are not declared solved by the linear/free-service theorem; nor are nonuniform tariffs claimed necessary for hardness in every larger subclass. These distinctions are explicit mathematical domains, not replacements for the original model.

\section*{2. Encoded language, parameters and the reduction}
Section~\ref{sec:decision60} defines the finite rational decision language, binary representation, threshold, validation and infeasible-input convention. The lower bound remains a parameterized many-one reduction for the original realized-ceiling model. We retain the group-sum construction, baseline-forcing argument, slot count, no-carry consistency checks and complete proof of Theorem~\ref{thm:wone59}. The text states why hardness in the sum also rules out fixed-parameter algorithms in either menu allowance or cap count alone. We make no W[1]-membership, completeness, strong-NP-hardness, exact-exponent, ETH or kernelization claim. Distinct fee values are not confused with exceptional commands. Table~\ref{tab:frontier60} and the expanded literature discussion compare the result with parameterized knapsack work.

\section*{3. Full-graph versus length-bounded conservation}
\textbf{Concern.} Full-graph capacity conservation can be stronger than needed under an edge allowance; the recognition proof, witness and evaluation model need precise scope.

\textbf{Response.} Theorem~\ref{thm:recognition59} is retained with an explicit all-path qualifier. Theorem~\ref{thm:bounded60} in the companion supplies the exact length-bounded counterpart on a pruned layered graph. It recognizes conservation in $O(h(|V|+|E|))$ arithmetic operations and returns two admissible original paths with unequal capacities on failure. A three-edge example shows why full-graph recognition is sufficient but not necessary for the bounded problem. The proof treats multiple sources, elapsed edge count, zero-capacity terminal edges and same-path repair explicitly. Pointer-based witness construction avoids hiding quadratic path storage. The additive complexity does not acquire a second layer factor: the original dynamic program already tracks elapsed edges.

The companion also distinguishes an exact-value oracle model from polynomial bit evaluation of explicitly encoded rational piecewise-linear or piecewise-quadratic kernels. Numerical grid size is never replaced by its logarithm. New tests include conservative and nonconservative graphs, independently checked witnesses, bounded-path enumeration, zero-capacity and zero-deficit cases, exact repair and holey-support convolution.

\section*{4. A concrete interpretation beyond renewal commands}
Section~\ref{sec:corridor60} maps corridor rehabilitation with endogenous contract boundaries into the abstract resource-path model. An edge is an admissible contract, its capacity is physical corridor length, the allocation is treated length, its payoff is separable concave benefit net of mobilization, and the complement is untreated length. The potential and source-equality conditions follow from those units. The mapping states divisible-treatment and separability assumptions and does not assert validity for spillovers or indivisible work packages. This is a fully specified nonrenewal decision model, not field calibration or an empirical transportation claim.

\section*{5. Why price methods can outperform the deficit algorithm}
Section~\ref{sec:pricebridge60} adds Theorem~\ref{thm:pricebridge60}: equality of a fixed price bound with the original optimum holds exactly when one active book supports the accepted promise. A distance-to-support bound quantifies the price gap for an active book. Bracketing the promise with different books is not sufficient. The zero-fee linear case always has a one-price exact certificate. Conversely, an explicitly checked uniform-positive-fee example has a positive price gap even though the new tariff algorithm solves it in polynomial time. Thus tariff tractability is not being inferred from observed zero duality gaps.

The computational section retains the earlier adverse deficit results: 5 of 72 rationally certified targets, versus 71 for price search and 12 for enumeration; 55 deficit intervals check within budget and 17 requests exceed their parent deadline. Seven price requests branch and maximum depth is 16. We explain the numerical-grid and independent-checker costs and distinguish an a priori approximation theorem from an instance-specific price certificate. We neither hide failures nor describe an unimplemented hybrid as tested. A declared tariff applicability guard and the same-book support test provide concrete method-selection logic.

\section*{6. Experimental denominators, budgets and evidence classes}
The historical panel is correctly described as 56 distinct specifications, 72 tolerance-tagged cases and 288 requests. No execution is relabeled as a new independent instance. The new prospective design contains 90 distinct specifications, 270 budget-tagged cases and 1,260 requests, from three 30-seed families. Standard tariffs, two exceptional tariffs and curved/costly service isolate the theorem's applicability boundary. Every specification is attempted at 0.5, 1.5 and 3 seconds; the optimization cutoff is 60 percent of the total. Sources, all inputs and deterministic method orders are frozen and committed before timing. All completed, interrupted, failed and state-limited requests are retained. The new panel's realized counts are in Tables~\ref{tab:newrational60} and~\ref{tab:newnumerical60}; the companion and machine-readable summaries provide conditional phase-cost denominators and tails.

The families are model-derived service-contract experiments, not customer observations or a canonical external benchmark. Their joint generator and units are disclosed. Shared seeds, interventions and repeated time allowances make requests dependent, so their total is not used as an independent statistical sample size. The new evidence tests exact tariff compression and the stated short-budget certification question; it does not claim superiority on an unexecuted external problem class or unrestricted long-budget solver tuning.

\section*{7. Numerical solver semantics and proof-checking cost}
The primary tables separate rationally certified target attainment from numerical-bound target attainment. For price, deficit, enumeration and tariff, success requires an independent rational global interval and an original feasible policy completed within the parent allowance. For SCIP, the formulation is algebraically exact but the upper bound remains numerical. The selected book is reoptimized in rational arithmetic; this produces an exact feasible lower policy and can change the raw continuous solution. It does not certify the solver's numerical upper bound. The revised text states this projection convention explicitly.

Total time includes startup, setup, optimization, serialization, compression and checking. Per-family--budget--method summaries include medians, 90th and 95th percentiles, maxima, proof bytes, compressed bytes and phase denominators. Missing phases are not zeroed. Independent post-study verification does not retroactively change deadline success. Exact cap counts, joint types, weight-denominator bit lengths and predicted grid sizes accompany each input. The two evidence classes also have separate attainment curves.

\section*{8. Closest literature and prototype selection}
The introduction and Table~\ref{tab:frontier60} add modern bucket-graph labeling, heuristic-guided constrained-path search, fixed-charge nonlinear allocation via perspective Benders decomposition and parameterized knapsack comparisons. These methods are identified as antecedents, not claimed as directly benchmarked implementations. Published versions replace outdated preprint-only entries for menu-size approximation, correlated-goods complexity and dual quantization. The cap-preserving weighted-median grouping theorem is retained as a model-specific loss frontier; generic medoid and one-dimensional clustering facts are credited rather than advertised as new algorithms.

\section*{9. Preservation, standalone readers and reproducibility}
The complete previous main article and companion are retained byte-for-byte. All mathematical sections from the reviewed main manuscript and companion remain in the current readers, with a source-to-location preservation map. Historical revision directories and the review branch are unchanged. Current sources expand their scientific dependencies, and a flat code/data package includes a current runtime, original frozen execution snapshots, complete records and certificate checkers. Path-only runtime adaptations are mapped separately from the hashes of the code actually timed. No missing historical log is reconstructed by inventing a new execution.

The build records abstract length, input hashes, page counts, essential theorem labels, citation/reference diagnostics and PDF hashes. A read-only verification of the published checkout binds the final payload to its commit. An execution source freeze and a final publication commit are distinct records; the response does not conflate them.

\section*{10. Presentation and submission format}
The manuscript uses anonymous 11-point type, one-inch margins and one-and-a-half spacing. The abstract is text-only and below 200 words; the introduction contains no mathematical notation. References are author--year and alphabetized, and tables follow the references. All retained and new proofs are typeset rather than replaced by repository pointers. The submission checklist records the actual page category and companion-length comparison after compilation. The response supplies precise theorem and section cross-references rather than an acceptance prediction.

\paragraph{Remaining boundaries.} The revision establishes the stated new positive regime and retains the original lower bound and broader approximation results. It does not claim a field validation, a universal runtime advantage for deficit methods, or a uniform-fee theorem outside its proved linear/free-service domain. These limitations are kept visible so that the referee can evaluate the strengthened results on their actual assumptions.
